\documentclass[10pt,a4paper]{article}
\usepackage[margin=2.0cm,top=2.0cm,bottom=2.2cm]{geometry}
\usepackage{fontspec}
\usepackage{amsmath,amssymb}
\usepackage{booktabs}
\usepackage{tabularx}
\usepackage{fancyvrb}
\usepackage{xcolor}
\usepackage{enumitem}
\usepackage{fancyhdr}
\usepackage[colorlinks=true,linkcolor=blue!45!black,urlcolor=blue!45!black]{hyperref}

\setmainfont{DejaVu Sans}[Scale=0.88]
\setsansfont{DejaVu Sans}[Scale=0.88]
\setmonofont{DejaVu Sans Mono}[Scale=0.82]

\definecolor{rulegray}{gray}{0.72}
\newenvironment{quoting}%
  {\begin{list}{}{\setlength{\leftmargin}{1.4em}\setlength{\rightmargin}{1.4em}%
   \setlength{\topsep}{5pt}}\item[]\itshape}%
  {\end{list}}
\setlength{\parindent}{0pt}
\setlength{\parskip}{0.55em}
\renewcommand{\arraystretch}{1.18}

\pagestyle{fancy}\fancyhf{}
\fancyhead[L]{\footnotesize\sffamily Sentinel-1 input to the decoder}
\fancyhead[R]{\footnotesize\sffamily 2026-08-24}
\fancyfoot[C]{\footnotesize\thepage}
\renewcommand{\headrulewidth}{0.3pt}

\newcommand{\code}[1]{\texttt{\small #1}}
\DefineVerbatimEnvironment{blk}{Verbatim}%
  {fontsize=\scriptsize,frame=leftline,framerule=0.6pt,%
   rulecolor=\color{rulegray},xleftmargin=6pt,framesep=6pt}

\title{\bfseries\sffamily Computing the Sentinel-1 input\\to the decoder}
\author{\sffamily Soil moisture project}
\date{\sffamily 2026-08-24}

\begin{document}
\maketitle
\thispagestyle{fancy}
{\small\tableofcontents}
\medskip\hrule\medskip

\textbf{Scope.} This document defines one thing only: what is computed from the Sentinel-1
archive and handed to each stage of the decoder. It does not discuss targets, losses,
auxiliary heads, terrain or thermal data.

\medskip\hrule\medskip

\section{Why anything is needed}

The decoder turns a $14\times14$ grid into a $112\times112$ map --- 196 values per channel
becoming 12{,}544. Every one of those values must come from somewhere.

In a conventional U-Net it comes from the skip connections: the encoder passed through
$112\times112$ on the way down and saved a copy. TerraMind is a Vision Transformer. It cuts
the image into $16\times16$-pixel patches once, at the start, and every subsequent layer
operates on those 196 patches. It never visits $112\times112$, so there is no
$112\times112$ copy to save --- \code{L3}, \code{L6}, \code{L9} and \code{L12} are
\textbf{all} $14\times14$ (\code{model.py:530}).

At \code{model.py:253-255} those skips are FiLM-modulated; at \code{model.py:257-264} they
are interpolated up and concatenated. So every decoder stage receives 160\,m information
relabelled as 80/40/20\,m. \textbf{Between 160\,m and 20\,m the decoder has no measurement
at all}, and fills the gap from learned convolution weights. §23 measured what it fills it
with: confident structure up to 0.19\,m³/m³, strongest at the \emph{flattest} site.

Sentinel-1 is the only dense observation available that varies in \textbf{both} space and
time. It sits on disk as a $224\times224$ raster of $\sigma^0$ in dB per acquisition,
median 440 acquisitions per station.

\begin{blk}
   s1_asc   993 stations   median 212   p10  90   max 958
   s1_desc  910 stations   median 238   p10   5   max 956
   asc+desc                median 440   p10 144   only 9 stations below 50

   /projects/prjs1968/satellite_zarr/{station}.zarr/s1_{asc,desc}/
       data   (N, 2, 224, 224)  fp16, dB, bands [VV, VH]
       dates  (N,)              b'YYYYMMDD'
\end{blk}

\section{Notation}

\begin{center}\small
\begin{tabularx}{\linewidth}{lX}
\toprule
symbol & meaning \\
\midrule
$p$ & acquisition (pass) index, $p = 1 \dots n$ \\
$k$ & cell index at a given grid level \\
$\ell$ & decoder level, $\ell \in \{14,\,28,\,56,\,112\}$ \\
$f_\ell$ & pixels per cell, $f_\ell = 224/\ell$ \\
$B_k$ & the block of $f_\ell^2$ pixels belonging to cell $k$ \\
$P$ & backscatter in \textbf{linear power} \\
$T_\ell$ & backscatter in \textbf{dB} after multilooking to level $\ell$ \\
\bottomrule
\end{tabularx}
\end{center}

The four levels correspond to the four decoder stages:
\[
\ell = 14 \to 160\,\mathrm{m}, \quad
28 \to 80\,\mathrm{m}, \quad
56 \to 40\,\mathrm{m}, \quad
112 \to 20\,\mathrm{m}.
\]

\section{Multilooking, and why it must be done in linear power}

Radar backscatter is multiplicative and is delivered as a logarithm. To average $N$ pixels
correctly you must average the \emph{power}, then take the log:
\[
P(p,u,v) \;=\; 10^{\,\sigma_{dB}(p,u,v)/10}
\]
\[
T_\ell(p,k) \;=\; 10\log_{10}\!\left(\frac{1}{|B_k|}\sum_{(u,v)\in B_k} P(p,u,v)\right),
\qquad |B_k| = f_\ell^{\,2}.
\]

Averaging in dB instead computes $\overline{10\log_{10} P}$, the \textbf{geometric} mean of
power, which is a different quantity:
\[
10\log_{10}\!\big(\overline{P}\big) \;\neq\; \overline{10\log_{10}(P)} .
\]

Measured on \code{AmeriFlux\_CA-Cbo}, the discrepancy between the two is a mean bias of
\textbf{0.53\,dB}, a maximum of \textbf{7.5\,dB}, and $r = 0.952$. The error is largest in
heterogeneous cells --- exactly where the signal lives.

\begin{quoting}
Each level must be computed from $P$ independently. A coarse level cannot be obtained by
pooling a finer $T$, for the same reason.
\end{quoting}

\section{Grouping by relative orbit}

Backscatter depends on the incidence angle, which is fixed per relative orbit. Mixing
orbits injects viewing geometry into everything downstream. Sentinel-1's repeat cycle is
exactly 12 days, so the relative orbit is constant modulo 12:
\[
g(p) \;=\; \operatorname{ord}(t_p) \bmod 12 .
\]

Two keys are \textbf{merged} only when their long-run means agree at the within-group
reproducibility ceiling. That recovers Sentinel-1A and 1B flying the \emph{same} relative
orbit six days apart, while keeping genuinely different orbits separate. Measured on
\code{CA-Cbo}, where the mod-12 key gives $\{0{:}116,\ 1{:}107,\ 7{:}108\}$ passes:

\begin{blk}
   split-half r(c_k) WITHIN a group   (the ceiling)     0.9979  0.9981  0.9982

   r(c_k) BETWEEN groups   0 vs 1   0.9632   different relative orbit -> keep apart
                           0 vs 7   0.9552   different relative orbit -> keep apart
                           1 vs 7   0.9952   AT the ceiling: S1A/S1B, same orbit -> MERGE
\end{blk}

Groups with fewer than 20 passes are discarded (see §\ref{sec:npasses}).

\section{The decomposition}

Within one station, one orbit group and one polarisation, form the $n \times \ell^2$ matrix
$T_\ell$ and define three means:
\[
\mu \;=\; \frac{1}{n\,\ell^2}\sum_{p}\sum_{k} T_\ell(p,k),
\qquad
r_p \;=\; \frac{1}{\ell^2}\sum_{k} T_\ell(p,k),
\qquad
c_k \;=\; \frac{1}{n}\sum_{p} T_\ell(p,k),
\]
and the remainder
\[
\boxed{\; d(p,k) \;=\; T_\ell(p,k) \;-\; r_p \;-\; c_k \;+\; \mu \;}
\]

Equivalently, a double-centring projection:
\[
D \;=\; \left(I_n - \tfrac{1}{n}\mathbf{1}\mathbf{1}^{\!\top}\right) T_\ell
        \left(I_{\ell^2} - \tfrac{1}{\ell^2}\mathbf{1}\mathbf{1}^{\!\top}\right),
\]
from which the defining properties are immediate:
\[
\sum_k d(p,k) = 0 \;\; \forall p,
\qquad
\sum_p d(p,k) = 0 \;\; \forall k .
\]

The three components are orthogonal, so the variance splits exactly:
\[
\operatorname{Var}(T) \;=\;
\underbrace{\operatorname{Var}(r)}_{\text{season}} \;+\;
\underbrace{\operatorname{Var}(c)}_{\text{static map}} \;+\;
\underbrace{\operatorname{Var}(d)}_{\text{interaction}} .
\]

This is a check, not an assumption. Verified numerically to $2\times10^{-7}$\,dB²:

\begin{blk}
AmeriFlux_CA-Cbo        orbit 0    n=116
  Var(T) total                        =  2.1386 dB^2
  Var(r_p)  season                    =  0.4195    19.6 %
  Var(c_k)  static map                =  1.3618    63.7 %
  Var(d)    interaction               =  0.3573    16.7 %
  sum of the three                    =  2.1386    |sum - Var(T)| = 2.4e-07

ISMN_TxSON_CR1000-1     orbit 3    n=113     season 37.7 %  static 50.1 %  inter 12.2 %
ISMN_SCAN_Abrams        orbit 4    n=104     season 40.7 %  static 21.7 %  inter 37.7 %
\end{blk}

\section{What each component means physically}

\[
\sigma^0(p,k) \;=\;
\underbrace{c_k}_{\substack{\text{what this place \emph{is}}}} \;+\;
\underbrace{(r_p-\mu)}_{\substack{\text{how wet the}\\ \text{whole tile is}}} \;+\;
\underbrace{d(p,k)}_{\substack{\text{how wet \emph{this} spot is}\\ \text{versus its own norm}}}
\]

\begin{itemize}[leftmargin=1.2em,itemsep=2pt,topsep=3pt]
\item $c_k$ --- \textbf{permanent texture}. A wood is always bright, a smooth field always
dark. Roughness and structure, not water. It never changes.
\item $r_p - \mu$ --- \textbf{catchment wetness state} on date $p$. When it rains the
entire tile brightens together.
\item $d(p,k)$ --- \textbf{differential wetting and drying}. What is left once season and
permanent texture are removed.
\end{itemize}

A worked example on real \code{CA-Cbo} tokens (orbit 0, VV), three cells and two dates:

\begin{blk}
STEP 0  raw token sigma0 (dB)                     tile mean r_p
  20171217    -17.42    -9.70    -6.24       -10.05
  20200710     -6.85    -6.52    -3.47        -6.88
  grand mean mu = -8.40 dB
  per-location mean c_k =  -12.97    -8.15    -4.29

STEP 1  subtract the TILE MEAN FOR THAT DATE
  20171217     -7.37     0.35     3.81
  20200710      0.02     0.36     3.41

STEP 2  subtract the PER-LOCATION MEAN, add back mu       -> this is d
  20171217     -2.80     0.10    -0.30
  20200710      4.60     0.11    -0.70
\end{blk}

The left-hand cell swings 10.6\,dB in the raw data, but it is a permanently dark cell
($c_k = -12.97$), so most of that is simply what it is. After both subtractions it goes
$-2.80 \to +4.60$: drier than its own norm on one date, much wetter on the other. The
centre cell meanwhile barely moves, $+0.10 \to +0.11$ --- it just tracks the tile.

Note there is no separate ``seasonal mean'' to compute. $r_p$, the tile-wide mean
\emph{on that date}, \textbf{is} the season. The two subtractions commute.

\section{Why the decomposition rather than raw $\sigma^0$}

The three components sum back to raw $\sigma^0$ up to a constant, so feeding all three is
\textbf{lossless}: everything raw contains is present, merely factored.

The reason to factor it is that \textbf{$c_k$ cannot be derived from a single sample}. It
requires the multi-year archive. A network given raw $\sigma^0$ could only recover $c_k$ by
memorising a per-pixel mean for each of 993 stations --- which is precisely the
memorisation the design is trying to avoid, and which does not transfer to a station the
model has never seen.

Computing $c_k$ offline injects information the network cannot otherwise obtain, and it
\emph{does} transfer, because at inference $c_k$ is computed from the new station's own
archive.

\begin{quoting}
Worth testing raw $\sigma^0$ as an ablation, specifically on out-of-sample stations. The
prediction is that raw matches on train and collapses out of sample, because that is where
a memorised $c_k$ stops working.
\end{quoting}

\section{Normalisation}

Each quantity is standardised within its station--orbit group, so that no station dominates:
\[
\tilde d = \frac{d}{\operatorname{sd}(d)},
\qquad
\tilde c = \frac{c - \bar c}{\operatorname{sd}(c)},
\qquad
w_p = \frac{r_p - \bar r}{\operatorname{sd}(r)} .
\]

$w_p$ is the catchment wetness state as a single standardised scalar per date.

\begin{quoting}
This is unrelated to the global z-score TerraMind applies at \code{precompute\_terramind.py:140},
which uses one fixed constant (\code{S1RTC} mean $[-10.93, -17.329]$, sd $[4.391, 4.459]$)
for every pixel of every image. A global affine map cannot remove a term that varies by
date, nor one that varies by location. The two operations do not overlap.
\end{quoting}

\section{Missing cells}

With a validity mask $m(p,k) \in \{0,1\}$ the single-pass centring is no longer exact.
Alternate masked removals until both margins vanish:
\[
d^{(i+\frac12)} = d^{(i)} - \frac{\sum_k m\,d^{(i)}}{\sum_k m},
\qquad
d^{(i+1)} = d^{(i+\frac12)} - \frac{\sum_p m\,d^{(i+\frac12)}}{\sum_p m} .
\]

Terminate when $\max_p\big|\overline{d}(p,\cdot)\big| < 10^{-4}$ and
$\max_k\big|\overline{d}(\cdot,k)\big| < 10^{-4}$. Typically 3--5 iterations. Measured
missing-cell rates are small:

\begin{blk}
station                       N  cell invalid%  dates w/ any bad   orbits (n>=20)
AmeriFlux_CA-Cbo            332       0.000%              0.0%       [0, 1, 7]
ISMN_TxSON_CR1000-1         114       0.000%              0.0%             [3]
ISMN_SCAN_Abrams            275       0.000%              0.0%         [4, 10]
ISMN_SCAN_AdamsRanch#1      239       0.201%              0.4%          [0, 6]
\end{blk}

\section{How many passes $c_k$ needs}
\label{sec:npasses}

$c_k$ is a mean over $n$ passes, so it carries sampling error. Measured on \code{CA-Cbo}
orbit 0, where the spread of $c_k$ across cells is 1.167\,dB:

\begin{center}\small
\begin{tabular}{rrrr}
\toprule
passes used & $r(c_{\text{sub}}, c_{\text{full}})$ & RMSE (dB) & as \% of the $c_k$ spread \\
\midrule
 5 & 0.9791 & 0.351 & 30.1 \\
10 & 0.9897 & 0.257 & 22.0 \\
20 & 0.9954 & 0.169 & 14.5 \\
30 & 0.9972 & 0.137 & 11.7 \\
50 & 0.9988 & 0.089 &  7.6 \\
80 & 0.9996 & 0.051 &  4.3 \\
\bottomrule
\end{tabular}
\end{center}

$n \ge 50$ is comfortable; $n \ge 20$ is the floor. Medians are 212 (ascending) and 238
(descending), so almost every station qualifies. The degrees of freedom consumed by fitting
$r_p$ and $c_k$ are 5.5\,\% at $n=20$ and 1.5\,\% at $n=100$ --- negligible.

\begin{quoting}
$c_k$ should also be checked for seasonal balance. Pass counts are 22--28\,\% per season at
every station examined, so the column mean is a genuine climatology rather than a
summer-weighted one. A station with a lopsided record would give a $c_k$ that absorbs part
of the interaction signal.
\end{quoting}

\section{Note on grid alignment}

Within one tile, all cells share the same acquisition dates. One Sentinel-1 acquisition is a
single $224\times224$ scene covering the whole tile: \code{data} is \code{(N,2,224,224)} and
\code{dates} is \code{(N,)} --- one entry per \emph{image}, not per location. $T_\ell$ is
therefore a complete $n \times \ell^2$ matrix and $c_k$ is a plain column mean over it. No
temporal alignment step exists or is needed.

Dates do differ across stations, and across orbit groups within a station. Both are
harmless, because centring is always performed \emph{within} one station and one orbit
group; $c_k$ is never compared across stations, and $\tilde d$ is standardised per group.

\section{What is handed to each decoder stage}

Per sample, using the most recent available pass $p^\star$:

\begin{center}\small
\begin{tabularx}{\linewidth}{llrX}
\toprule
stage & grid & cell & channels \\
\midrule
bottleneck & $14\times14$   & 160\,m & $\tilde c_{14},\ \tilde d_{14}(p^\star),\ w_{p^\star}$ \\
\code{up1} & $28\times28$   &  80\,m & $\tilde c_{28},\ \tilde d_{28}(p^\star),\ w_{p^\star}$ \\
\code{up2} & $56\times56$   &  40\,m & $\tilde c_{56},\ \tilde d_{56}(p^\star),\ w_{p^\star}$ \\
\code{up3} & $112\times112$ &  20\,m & $\tilde c_{112},\ \tilde d_{112}(p^\star),\ w_{p^\star}$ \\
\bottomrule
\end{tabularx}
\end{center}

$w_{p^\star}$ is a scalar broadcast across the grid. Two housekeeping scalars accompany
them: the age of $p^\star$ in days, and a validity flag which is zero when no usable pass
exists --- in which case $\tilde d = 0$, the natural ``behaving like the tile average''
value.

Concatenation alongside the existing skip, at \code{model.py:257-264}:

\begin{blk}
x = self.bottle_proj(bottleneck)                                 # (B,512,14,14)
x = torch.cat([x, c_14, d_14, w], 1)                             #  512 + 3

x = self.up1(x)                                                  # 14 -> 28
x = self.conv1(torch.cat([x, film9(skip_L9),  c_28,  d_28,  w], 1))   # 512+512+3 -> 256
x = self.up2(x)                                                  # 28 -> 56
x = self.conv2(torch.cat([x, film6(skip_L6),  c_56,  d_56,  w], 1))   # 256+256+3 -> 128
x = self.up3(x)                                                  # 56 -> 112
x = self.conv3(torch.cat([x, film3(skip_L3),  c_112, d_112, w], 1))   # 128+128+3 ->  64
\end{blk}

Each \code{in\_channels} grows by 3; \code{\_ConvBlock} (\code{model.py:171}) is otherwise
untouched. Added parameters are $3 \times C_{\text{out}} \times 9$ per stage:
$6912 + 3456 + 1728 \approx 12$\,k.

The weights on the new channels are zero-initialised,
\begin{blk}
for c in (conv1, conv2, conv3):
    c.block[0].weight[:, -3:] = 0
\end{blk}
so at initialisation they contribute exactly nothing, the forward pass is bit-identical to
\code{best.pt}, and §26's provenance gate applies unchanged. Training moves them off zero.

\section{Storage and build cost}

Per station $\times$ orbit group $\times$ polarisation, with $n \approx 200$ passes:
\[
\big(14^2 + 28^2 + 56^2 + 112^2\big) \times n
\;=\; 16{,}660 \times 200 \;\approx\; 3.3\times10^{6}\ \text{values},
\]
which at fp16 is 6.6\,MB, so roughly \textbf{13\,GB} across 993 stations. The $c_\ell$
arrays are one per level and negligible.

Build cost is dominated by reading the rasters: $993 \times 2 \times {\sim}66$\,MB
$\approx 130$\,GB. One \code{sbatch} with \code{Pool(64)} and
\code{--cpus-per-task=64}, approximately 30--60 minutes.

\section{Summary}

\begin{blk}
   RAW           sigma0(p,k)  in dB, (N,2,224,224) per station per orbit direction
      |
      |  P = 10^(sigma/10)                       to LINEAR power
      |  mean over each f x f block, back to dB  MULTILOOK to level l
      v
   T_l(p,k)      (n, l^2) in dB
      |
      |  group by  ord(t_p) mod 12               RELATIVE ORBIT
      |  merge keys at the split-half ceiling
      |  drop groups with n < 20
      v
   per group:  mu, r_p, c_k
      |
      |  d = T - r_p - c_k + mu                  DOUBLE-CENTRE
      |  iterate if cells are missing
      |  standardise
      v
   FED TO STAGE l:   c~_l   (permanent texture,      per cell)
                     d~_l   (today's anomaly,        per cell)
                     w_p    (catchment wetness state, scalar)
\end{blk}

\end{document}
